\documentclass[12pt,a4paper]{ctexbook}
\usepackage{geometry}
\geometry{margin=2.2cm}
\usepackage{setspace}
\setstretch{1.35}
\usepackage{hyperref}
\title{全唐诗选·poet.tang.1000}
\author{古典文献整理}
\date{}
\begin{document}
\maketitle
\tableofcontents
\newpage
\section{橫吹曲辭 洛陽陌}
\textbf{李白}\par\vspace{0.5em}
白玉誰家郎，回車渡天津。\par
看花東上陌，驚動洛陽人。\par

\section{橫吹曲辭 長安道}
\textbf{崔顥}\par\vspace{0.5em}
長安甲第高入雲，誰家居住霍將軍。\par
日晚朝回擁賓從，路傍拜揖何紛紛。\par
莫言炙手手可熱，須臾火盡灰亦滅。\par
莫言貧賤即可欺，人生富貴自有時。\par
一朝天子賜顏色，世上悠悠應始知。\par

\section{橫吹曲辭 長安道}
\textbf{孟郊}\par\vspace{0.5em}
胡風激秦樹，賤子風中泣。\par
家家朱門開，得見不可入。\par
長安十二衢，投樹鳥亦急。\par
高閣何人家？笙簧正喧吸。\par

\section{橫吹曲辭 長安道}
\textbf{顧況}\par\vspace{0.5em}
長安道，人無衣，馬無草，何不歸來山中老。\par

\section{橫吹曲辭 長安道}
\textbf{聶夷中}\par\vspace{0.5em}
此地無駐馬，夜中猶走輪。\par
所以路旁草，少於衣上塵。\par

\section{橫吹曲辭 長安道}
\textbf{韋應物}\par\vspace{0.5em}
漢家宮殿含雲煙，兩宮十里相連延。\par
晨霞出沒弄丹闕，春雨依微自甘泉。\par
春雨依微春尚早，長安貴遊愛芳草。\par
寶馬橫來下建章，香車却轉避馳道。\par
貴遊誰最貴，衛霍世難比。\par
何能蒙主恩？幸遇邊塵起。\par
歸來甲第拱皇居，朱門峨峨臨九衢。\par
中有流蘇合歡之寶帳，一百二十鳳皇羅列含明珠。\par
下有錦鋪翠被之粲爛，博山吐香五雲散。\par
麗人綺閣情飄颻，頭上鴛釵雙翠翹。\par
低鬟曳袖回春雪，聚黛一聲愁碧霄。\par
山珍海錯棄藩籬，烹犢炰羔如折葵。\par
既請列侯封部曲，還將金印授廬兒。\par
歡榮若此何所苦，但苦白日西南馳。\par

\section{橫吹曲辭 長安道}
\textbf{白居易}\par\vspace{0.5em}
花枝缺處青樓開，豔歌一曲酒一杯。\par
美人勸我急行樂，自古朱顏不再來。\par
君不見外州客，長安道，一回來，一回老。\par

\section{橫吹曲辭 長安道}
\textbf{薛能}\par\vspace{0.5em}
汲汲復營營，東西連兩京。\par
關繻古若在，山岳累應成。\par
各自有身事，不相知姓名。\par
交馳喧衆類，分散入重城。\par
此路去無盡，萬方人始生。\par
空餘片言苦，來往覓劉楨。\par

\section{橫吹曲辭 長安道}
\textbf{僧貫休}\par\vspace{0.5em}
憧憧合合，八表一轍。\par
黃塵霧合，車馬火熱。\par
名湯風雨，利輾霜雪。\par
千車萬駄，半宿關月。\par
上有堯禹，下有夔契。\par
紫氣銀輪兮常覆金闕，仙掌捧日兮濁河澄澈。\par
愚將草木兮有言，與華封人兮不別。\par

\section{橫吹曲辭 長安道}
\textbf{沈佺期}\par\vspace{0.5em}
秦地平如掌，層城出雲漢。\par
樓閣九衢春，車馬千門旦。\par
綠柳開復合，紅塵聚還散。\par
日晚鬬雞回，經過狹斜看。\par

\section{橫吹曲辭 梅花落}
\textbf{盧照鄰}\par\vspace{0.5em}
梅嶺花初發，天山雪未開。\par
雪處疑花滿，花邊似雪回。\par
因風入舞袖，雜粉向妝臺。\par
匈奴幾萬里，春至不知來。\par

\section{橫吹曲辭 梅花落}
\textbf{沈佺期}\par\vspace{0.5em}
鐵騎幾時回，金閨怨早梅。\par
雪中花已落，風暖葉應開。\par
夕逐新春管，香迎小歲杯。\par
感時何足貴，書裏報輪臺。\par

\section{橫吹曲辭 梅花落}
\textbf{劉方平}\par\vspace{0.5em}
新歲芳梅樹，繁苞四面同。\par
春風吹漸落，一夜幾枝空。\par
小婦今如此，長城恨不窮。\par
莫將遼海雪，來比後庭中。\par

\section{橫吹曲辭 紫騮馬}
\textbf{盧照鄰}\par\vspace{0.5em}
騮馬照金鞍，轉戰入皐蘭。\par
塞門風稍急，長城水正寒。\par
雪暗鳴珂重，山長噴玉難。\par
不辭橫絕漠，流血幾時乾。\par

\section{橫吹曲辭 紫騮馬}
\textbf{李白}\par\vspace{0.5em}
紫騮行且嘶，雙飜碧玉蹄。\par
臨流不肎渡，似惜錦障泥。\par
白雪關山遠，黃雲海樹迷。\par
揮鞭萬里去，安得念春閨。\par

\section{橫吹曲辭 紫騮馬}
\textbf{李益}\par\vspace{0.5em}
爭場看鬬雞，白鼻紫騧嘶。\par
漳水春歸晚，叢臺日向低。\par
歇鞍珠作汗，試劒玉如泥。\par
爲謝紅梁燕，年年妾獨棲。\par

\section{橫吹曲辭 紫騮馬}
\textbf{秦韜玉}\par\vspace{0.5em}
渥洼奇骨本難求，況是豪家重紫騮。\par
臕大宜懸銀壓銙，力渾欺却玉銜頭。\par
生獰弄影風隨起，躞蹀衝塵汗滿溝。\par
若遇丈夫皆調御，任從騎取覓封侯。\par

\section{橫吹曲辭 驄馬}
\textbf{李羣玉}\par\vspace{0.5em}
浮雲何權奇，絕足勢未知。\par
長嘶青海風，躞蹀振雲絲。\par
由來渥洼種，本是蒼龍兒。\par
穆滿不再活，無人崑閬騎。\par
君識躍嶠怯，寧勞耀金羈。\par
青芻與白水，空笑駑駘肥。\par
伯樂儻一見，應驚耳長垂。\par
當思八荒外，逐日向瑤池。\par

\section{橫吹曲辭 驄馬曲}
\textbf{紀唐天}\par\vspace{0.5em}
連錢出塞蹋沙蓬，豈比當時御史驄。\par
逐北自諳深磧路，連嘶誰念靜邊功。\par
登山每與青雲合，弄影應知碧草同。\par
今日虜平將換妾，不如羅袖舞春風。\par

\section{橫吹曲辭 雨雪曲}
\textbf{李端}\par\vspace{0.5em}
天山一丈雪，雜雨夜霏霏。\par
溼馬胡歌亂，經烽漢火微。\par
丁零蘇武別，疎勒范羌歸。\par
若著關頭過，長榆葉定稀。\par

\section{橫吹曲辭 雨雪曲}
\textbf{翁綬}\par\vspace{0.5em}
邊聲四合殷河流，雨雪飛來徧隴頭。\par
鐵嶺探人迷鳥道，陰山飛將溼貂裘。\par
斜飄旌斾過戎帳，半雜風沙入戍樓。\par

\section{橫吹曲辭 劉生}
\textbf{盧照鄰}\par\vspace{0.5em}
劉生氣不平，抱劒欲專征。\par
報恩爲豪俠，死難在橫行。\par
翠羽裝劒鞘，黃金飾馬纓。\par
但令一顧重，不吝百身輕。\par

\section{橫吹曲辭 雍臺歌}
\textbf{溫庭筠}\par\vspace{0.5em}
太子池南樓百尺，八窗新樹疎簾隔。\par
黃金鋪首畫鉤陳，羽葆亭童拂交戟。\par
盤紆闌楯臨高臺，帳殿臨流鸞扇開。\par
早雁聲鳴細波起，映花鹵簿龍飛回。\par

\section{橫吹曲辭 捉搦歌}
\textbf{張祜}\par\vspace{0.5em}
門上關，牆上棘，窗中女子聲唧唧，洛陽大道徒自直。\par
女子心在婆舍側，嗚嗚籠鳥觸四隅。\par
養男男娶婦，養女女嫁夫。\par
阿婆六十翁七十，不知女子長日泣，從他嫁去無悒悒。\par

\section{橫吹曲辭 幽州胡馬客歌}
\textbf{李白}\par\vspace{0.5em}
幽州胡馬客，綠眼虎皮冠。\par
笑拂兩隻箭，萬人不可干。\par
彎弓若轉月，白雁落雲端。\par
雙雙掉鞭行，遊獵向樓蘭。\par
出門不顧後，報國死何難。\par
天驕五單于，狼戾好凶殘。\par
牛馬散北海，割鮮若虎餐。\par
雖居燕支山，不道朔雪寒。\par
婦女馬上笑，顏如赬玉盤。\par
飜入射鳥獸，花月醉雕鞍。\par
旄頭四光芒，爭戰若蜂攢。\par
白刃灑赤血，流沙爲之丹。\par
名將古誰是，疲兵良可歎。\par
何時天狼滅？父子得閑安。\par

\section{橫吹曲辭 白鼻騧}
\textbf{李白}\par\vspace{0.5em}
銀鞍白鼻騧，綠地障泥錦。\par
細雨春風花落時，揮鞭且就胡姬飲。\par

\section{橫吹曲辭 白鼻騧}
\textbf{張祜}\par\vspace{0.5em}
爲底胡姬酒，長來白鼻騧。\par
摘蓮拋水上，郎意在浮花。\par

\section{相和歌辭 箜篌引}
\textbf{李賀}\par\vspace{0.5em}
公乎公乎，提壺將焉如？\par
屈平沈湘不足慕，徐衍入海誠爲愚。\par
公乎公乎，牀有菅席盤有魚。\par
北里有賢兄，東鄰有小姑。\par
隴畝油油黍與葫，瓦甒濁醪蟻浮浮。\par
黍可食，醪可飲，公乎公乎其奈居。\par
被髮奔流竟何如？賢兄小姑哭嗚嗚。\par

\section{相和歌辭 公無渡河}
\textbf{李白}\par\vspace{0.5em}
黃河西來決崐崘，咆吼萬里觸龍門。\par
波滔天，堯咨嗟。\par
大禹理百川，兒啼不窺家。\par
殺湍湮洪水，九州始蠶麻。\par
其害乃去，茫然風沙。\par
被髮之叟狂而癡，清晨徑流欲奚爲。\par
旁人不惜妻止之，公無渡河苦渡之。\par
虎可搏，河難憑。\par
公果溺死流海湄，有長鯨白齒若雪山。\par
公乎公乎，挂骨於其間，箜篌所悲竟不還。\par

\section{相和歌辭 公無渡河}
\textbf{王建}\par\vspace{0.5em}
渡頭惡天兩岸遠，波濤塞川如疊坂。\par
幸無白刃驅向前，何用將身自棄捐。\par
蛟龍齧屍魚食血，黃泥直下無青天。\par
男兒縱輕婦人語，惜君性命還須取。\par
婦人無力挽斷衣，舟沈身死悔難追，公無渡河公自爲。\par

\section{相和歌辭 公無渡河}
\textbf{溫庭筠}\par\vspace{0.5em}
黃河怒浪連天來，大響谹谹如殷雷。\par
龍伯驅風不敢上，百川噴雪高崔嵬。\par
二十五弦何太哀？請公勿渡立裴回。\par
下有狂蛟鋸爲尾，裂帆截櫂磨霜齒。\par
神錐鑿石塞神潭，白馬䟃𧽼赤塵起。\par
公乎躍馬揚玉鞭，滅沒高蹄日千里。\par

\section{相和歌辭 公無渡河}
\textbf{王叡}\par\vspace{0.5em}
濁波洋洋兮凝曉霧，公無渡河兮公苦渡。\par
風號水激兮呼不聞，提壺看入兮中流去。\par
浪擺衣裳兮隨步沒，沈屍深入兮蛟螭窟。\par
蛟螭盡醉兮君血乾，推出黃沙兮泛君骨。\par
當時君死妾何適？遂就波瀾合魂魄。\par
願持精衛銜石心，窮取河源塞泉脈。\par

\section{相和歌辭 江南曲}
\textbf{宋之問}\par\vspace{0.5em}
妾住越城南，離居不自堪。\par
採花驚曙鳥，摘葉餧春蠶。\par
嬾結茱萸帶，愁安玳瑁簪。\par
侍臣消瘦盡，日暮碧江潭。\par

\section{相和歌辭 江南曲}
\textbf{劉眘虛}\par\vspace{0.5em}
美人何蕩漾，湖上風月長。\par
玉手欲有贈，裴回雙鳴璫。\par
歌聲隨淥水，怨色起朝陽。\par
日暮還家望，雲波橫洞房。\par

\section{相和歌辭 江南曲}
\textbf{丁仙芝}\par\vspace{0.5em}
長干斜路北，近浦是兒家。\par
有意來相訪，明朝出浣紗。\par
發向橫塘口，船開值急流。\par
知郎舊時意，且請攏船頭。\par
昨暝逗南陵，風聲波浪阻。\par
入浦不逢人，歸家誰信汝。\par
未曉已成妝，乘潮去茫茫。\par
因從京口渡，使報邵陵王。\par
始下芙蓉樓，言發琅邪岸。\par
急爲打船開，惡許傍人見。\par

\section{相和歌辭 江南曲八首 一}
\textbf{劉希夷}\par\vspace{0.5em}
暮宿南洲草，晨行北岸林。\par
日懸滄海闊，水隔洞庭深。\par
煙景無留意，風波有異潯。\par
歲遊難極目，春戲易爲心。\par
朝夕無榮遇，芳菲已滿襟。\par

\section{相和歌辭 江南曲八首 二}
\textbf{劉希夷}\par\vspace{0.5em}
豔唱潮初落，江花露未晞。\par
春洲驚翡翠，朱服弄芳菲。\par
畫舫煙中淺，青陽日際微。\par
錦帆衝浪溼，羅袖拂行衣。\par
含情罷所采，相歎惜流暉。\par

\section{相和歌辭 江南曲八首 三}
\textbf{劉希夷}\par\vspace{0.5em}
君爲隴西客，妾遇江南春。\par
朝遊含靈果，夕采弄風蘋。\par
果氣時不歇，蘋花日自新。\par
以此江南物，持贈隴西人。\par
空盈萬里懷，欲贈竟無因。\par

\section{相和歌辭 江南曲八首 四}
\textbf{劉希夷}\par\vspace{0.5em}
皓如楚江月，靄若吳岫雲。\par
波中自皎鏡，山上亦氛氳。\par
明月留照妾，輕雲持贈君。\par
山川各離散，光氣乃殊分。\par
天涯一爲別，江北自相聞。\par

\section{相和歌辭 江南曲八首 五}
\textbf{劉希夷}\par\vspace{0.5em}
艤舟乘潮去，風帆振草涼。\par
潮平見楚甸，天際望維揚。\par
洄泝經千里，煙波接兩鄉。\par
雲明江嶼出，日照海流長。\par
此中逢歲晏，浦樹落花芳。\par

\section{相和歌辭 江南曲八首 六}
\textbf{劉希夷}\par\vspace{0.5em}
暮春三月晴，維揚吳楚城。\par
城臨大江汜，迴映洞浦清。\par
晴雲曲金閣，珠樓碧煙裏。\par
月明芳樹羣鳥飛，風過長林雜花起。\par
可憐離別誰家子，於此一至情何已。\par

\section{相和歌辭 江南曲八首 七}
\textbf{劉希夷}\par\vspace{0.5em}
北堂紅草盛丰茸，南湖碧水照芙蓉。\par
朝遊暮起金花盡，漸覺羅裳珠露濃。\par
自惜妍華三五歲，已歎關山千萬重。\par
人情一去無還日，欲贈懷芳怨不逢。\par

\section{相和歌辭 江南曲八首 八}
\textbf{劉希夷}\par\vspace{0.5em}
憶昔江南年盛時，平生怨在長洲曲。\par
冠蓋星繁湘水上，衝風摽落洞庭淥。\par
落花舞袖紅紛紛，朝霞高閣洗晴雲。\par
誰言此處嬋娟子，珠玉爲心以奉君。\par

\section{相和歌辭 江南曲}
\textbf{于鵠}\par\vspace{0.5em}
偶向江邊采白蘋，還隨女伴賽江神。\par
衆中不敢分明語，暗擲金錢卜遠人。\par

\section{相和歌辭 江南曲}
\textbf{李益}\par\vspace{0.5em}
嫁得瞿塘賈，朝朝誤妾期。\par
早知潮有信，嫁與弄潮兒。\par

\section{相和歌辭 江南曲}
\textbf{李賀}\par\vspace{0.5em}
汀洲白蘋草，柳惲乘馬歸。\par
江頭樝樹香，岸上蝴蝶飛。\par
酒杯若葉露，玉軫蜀桐虛。\par
朱樓通水陌，沙暖一雙魚。\par

\section{相和歌辭 江南曲}
\textbf{李商隱}\par\vspace{0.5em}
郎船安兩槳，儂舸動雙橈。\par
埽黛開宮額，裁裙約楚腰。\par
乖期方積思，臨醉欲拚嬌。\par
莫以采菱唱，欲羨秦臺簫。\par

\section{相和歌辭 江南曲}
\textbf{韓翃}\par\vspace{0.5em}
長樂花枝雨點消，江城日暮好相邀。\par
春樓不閉葳蕤鎖，綠水回通宛轉橋。\par

\section{相和歌辭 江南曲}
\textbf{溫庭筠}\par\vspace{0.5em}
妾家白蘋浦，日上芙蓉檝。\par
軋軋搖槳聲，移舟入茭葉。\par
溪長茭葉深，作底難相尋。\par
避郎郎不見，鸂鶒自浮沈。\par
拾萍萍無根，采蓮蓮有子。\par
不作浮萍生，寧作藕花死。\par
岸傍騎馬郎，烏帽紫游韁。\par
含愁復含笑，回首問橫塘。\par
妾住金陵步，門前朱雀航。\par
流蘇持作帳，芙蓉持作梁。\par
出入金犢幰，兄弟侍中郎。\par
前年學歌舞，定得郎相許。\par
連娟眉繞山，依約腰如杵。\par
鳳管悲若咽，鸞弦嬌欲語。\par
扇薄露紅鈆，羅輕壓金縷。\par
明月西南樓，珠簾玳瑁鉤。\par
橫波巧能笑，彎蛾不識愁。\par
花開子留樹，草長根依土。\par
早聞金溝遠，底事歸郎許。\par
不學楊白花，朝朝淚如雨。\par

\section{相和歌辭 江南曲}
\textbf{張籍}\par\vspace{0.5em}
江南人家多橘樹，吳姬舟上織白紵。\par
土地卑溼饒蟲蛇，連木爲牌入江住。\par
江村亥日長爲市，落帆渡橋來浦裏。\par
青莎覆城竹爲屋，無井家家飲潮水。\par
長江午日酤春酒，高高酒旗懸江口。\par
倡樓兩岸懸水柵，夜唱竹枝留北客。\par
江南風土歡樂多，悠悠處處盡經過。\par

\section{相和歌辭 江南曲}
\textbf{羅隱}\par\vspace{0.5em}
江煙溼雨鮫綃輭，漠漠遠山眉黛淺。\par
水國多愁又有情，夜槽壓酒銀船滿。\par
綳絲採怨凝曉空，吳王臺榭春夢中。\par
鴛鴦鸂鶒喚不起，平鋪淥水眠東風。\par
西陵路邊月悄悄，油壁輕車嫁蘇小。\par

\section{相和歌辭 江南曲}
\textbf{陸龜蒙}\par\vspace{0.5em}
爲愛江南雪，涉江聊采蘋。\par
水深煙浩浩，空對雙車輪。\par
車輪明月團，車蓋浮雲盤。\par
雲月徒自好，水中行路難。\par
遙遙洛陽道，夾岸生春草。\par
寄語櫂船郎，莫誇風浪好。\par

\section{相和歌辭 江南曲 一}
\textbf{陸龜蒙}\par\vspace{0.5em}
魚戲蓮葉間，參差隱葉扇。\par
鵁鶄鸀鳿窺，瀲灩無因見。\par

\section{相和歌辭 江南曲 二}
\textbf{陸龜蒙}\par\vspace{0.5em}
魚戲蓮葉東，初霞射紅尾。\par
傍臨謝山側，恰值清風起。\par

\section{相和歌辭 江南曲 三}
\textbf{陸龜蒙}\par\vspace{0.5em}
魚戲蓮葉西，盤盤舞波急。\par
潛依曲岸涼，正對斜光入。\par

\section{相和歌辭 江南曲 四}
\textbf{陸龜蒙}\par\vspace{0.5em}
魚戲蓮葉南，欹危午煙疊。\par
光搖越鳥巢，影亂吳娃檝。\par

\section{相和歌辭 江南曲 五}
\textbf{陸龜蒙}\par\vspace{0.5em}
魚戲蓮葉北，澄陽動微漣。\par
回看帝子渚，稍背鄂君船。\par

\section{相和歌辭 度關山}
\textbf{李端}\par\vspace{0.5em}
雁塞日初晴，胡關雪復平。\par
危竿緣廣漠，古竇傍長城。\par
拔劒金星出，彎弧玉羽鳴。\par
誰知係虜者，賈誼是書生。\par

\section{相和歌辭 關山曲 一}
\textbf{馬戴}\par\vspace{0.5em}
金鎖耀兜鍪，黃雲拂紫騮。\par
叛羌旗下戮，陷壁夜中收。\par
霜霰戎衣故，關河磧氣秋。\par
箭創殊未合，更遣擊蘭州。\par

\section{相和歌辭 關山曲 二}
\textbf{馬戴}\par\vspace{0.5em}
火發龍山北，中宵易左賢。\par
勒兵臨漢水，驚雁散胡天。\par
木落防河急，軍孤受敵偏。\par
猶聞漢皇怒，按劒待開邊。\par

\section{相和歌辭 登高丘而望遠}
\textbf{李白}\par\vspace{0.5em}
登高丘而望遠海，六鼇骨已霜，三山流安在。\par
扶桑半摧折，白日沈光彩。\par
銀臺金闕如夢中，秦皇漢武空相待。\par
精衛費木石，黿鼉無所憑。\par
君不見驪山茂陵盡灰滅，牧羊之子來攀登。\par
盜賊劫寶玉，精靈竟何能。\par

\section{相和歌辭 蒿里}
\textbf{僧貫休}\par\vspace{0.5em}
兔不遲，烏更急，但恐穆王八駿，著鞭不及。\par
所以蒿里，墳出蕺蕺，氣凌雲天，龍騰鳳集，盡爲風消土喫，狐掇蟻拾。\par
黃金不啼玉不泣，白楊騷屑。\par
亂風愁月，折碑石人。\par
莽穢榛沒，牛羊窸窣。\par
時見牧童兒，弄枯骨。\par

\section{相和歌辭 輓歌}
\textbf{趙微明}\par\vspace{0.5em}
寒日蒿上明，淒淒郭東路。\par
素車誰家子，丹旐引將去。\par
原下荆棘叢，叢邊有新墓。\par
人間痛傷別，此是長別處。\par
曠野何蕭條，青松白楊樹。\par

\section{相和歌辭 輓歌二首 一}
\textbf{于鵠}\par\vspace{0.5em}
陰風吹黃蒿，輓歌渡秋水。\par
車馬却歸城，孤墳月明裏。\par

\section{相和歌辭 輓歌二首 二}
\textbf{于鵠}\par\vspace{0.5em}
雙轍出郭門，緜緜東西道。\par
送死多於生，幾人得終老。\par
見人切肺肝，不如歸山好。\par
不聞哀哭聲，默默安懷抱。\par
時盡從物化，又免生憂擾。\par
世間壽者稀，盡爲悲傷惱。\par

\section{相和歌辭 輓歌}
\textbf{孟雲卿}\par\vspace{0.5em}
草草門巷喧，塗車儼成位。\par
冥寞何所須，盡我生人意。\par
北邙路非遠，此別終天地。\par
臨穴頻撫棺，至哀反無淚。\par
爾形未衰老，爾息猶童穉。\par
骨肉不可離，皇天若容易。\par
房帷即虛張，庭宇爲哀次。\par
薤露歌若斯，人生盡如寄。\par

\section{相和歌辭 輓歌}
\textbf{白居易}\par\vspace{0.5em}
丹旐何飛揚，素驂亦悲鳴。\par
晨光照閭巷，輀車儼欲行。\par
蕭條九月天，哀輓出重城。\par
借問送者誰，妻子與弟兄。\par
蒼蒼上古原，峨峨開新塋。\par
含酸一慟哭，異口同哀聲。\par
舊壟轉蕪絕，新墳日羅列。\par
春風草綠北邙山，此地年年生死別。\par

\section{相和歌辭 對酒}
\textbf{崔國輔}\par\vspace{0.5em}
行行日將夕，荒村古冢無人跡。\par
朦朧荆棘一鳥飛，屢唱提壺酤酒喫。\par
古人不達酒不足，遺恨精靈傳此曲。\par
寄言當代諸少年，平生且盡杯中淥。\par

\section{相和歌辭 對酒二首 一}
\textbf{李白}\par\vspace{0.5em}
松子棲金華，安期入蓬海。\par
此人古之仙，羽化竟何在？\par
浮生速流電，倏忽變光彩。\par
天地無凋換，容顏有遷改。\par
對酒不肯飲，含情欲誰待。\par

\section{相和歌辭 對酒二首 二}
\textbf{李白}\par\vspace{0.5em}
勸君莫拒杯，春風笑人來。\par
桃李如舊識，傾花向我開。\par
流鶯啼碧樹，明月窺金罍。\par
昨來朱顏子，今日白髮催。\par
棘生石虎殿，鹿走姑蘇臺。\par
自古帝王宅，城闕閉黃埃。\par
君若不飲酒，昔人安在哉。\par

\section{相和歌辭 陌上桑}
\textbf{李白}\par\vspace{0.5em}
美女渭橋東，春還事蠶作。\par
五馬如飛龍，青絲結金絡。\par
不知誰家子，調笑來相謔。\par
妾本秦羅敷，玉顏豔名都。\par
綠條映素手，采桑向城隅。\par
使君且不顧，況復論秋胡。\par
寒螿愛碧草，鳴鳳棲青梧。\par
託心自有處，但怪傍人愚。\par
徒令白日暮，高駕空踟躕。\par

\section{相和歌辭 陌上桑}
\textbf{常建}\par\vspace{0.5em}
翳翳陌上桑，南枝交北堂。\par
美人金梯出，素手自提筐。\par
非但畏蠶飢，盈盈嬌路傍。\par

\section{相和歌辭 陌上桑}
\textbf{陸龜蒙}\par\vspace{0.5em}
皓齒還如貝色含，長眉亦似煙華貼。\par
鄰娃盡著繡襠襦，獨自提筐采蠶葉。\par

\section{相和歌辭 采桑}
\textbf{郎大家宋氏}\par\vspace{0.5em}
春來南雁歸，日去西蠶遠。\par
妾思紛何極，客遊殊未返。\par

\section{相和歌辭 采桑}
\textbf{劉希夷}\par\vspace{0.5em}
楊柳送行人，青青西入秦。\par
秦家采桑女，樓上不勝春。\par
盈盈灞水曲，步步春芳綠。\par
紅臉耀明珠，絳脣含白玉。\par
回首渭橋東，遙憐樹色同。\par
青絲嬌落日，緗綺弄春風。\par
攜籠長歎息，逶迤戀春色。\par
看花若有情，倚樹疑無力。\par
薄暮思悠悠，使君南陌頭。\par
相逢不相識，歸去夢青樓。\par

\section{相和歌辭 采桑}
\textbf{李彥暐}\par\vspace{0.5em}
采桑畏日高，不待春眠足。\par
攀條有餘愁，那矜貌如玉。\par
千金豈不贈，五馬空踟躅。\par
何以變真性，幽篁雪中綠。\par

\section{相和歌辭 采桑}
\textbf{王建}\par\vspace{0.5em}
鳥鳴桑葉間，葉綠條復柔。\par
攀看去手近，放下長長鉤。\par
黃花蓋野田，白馬少年遊。\par
所念豈回顧，良人在高樓。\par

\section{相和歌辭 日出行}
\textbf{李白}\par\vspace{0.5em}
日出東方隈，似從地底來。\par
歷天又入海，六龍所舍安在哉？\par
其始與終古不息，人非元氣安能與之久裴回。\par
草不謝榮於春風，木不怨落於秋天。\par
誰揮鞭策驅四運，萬物興歇皆自然。\par
羲和羲和，汝奚汩沒於荒淫之波。\par
魯陽何德？駐景揮戈。\par
逆道違天，矯誣實多。\par
吾將囊括大塊，浩然與溟涬同科。\par

\section{相和歌辭 日出行}
\textbf{李賀}\par\vspace{0.5em}
白日下崐崘，發光如舒絲。\par
徒照葵藿心，不照遊子悲。\par
折折黃河曲，日從中央轉。\par
暘谷耳曾聞，若木眼不見。\par
奈何鑠石，胡爲銷人。\par
羿彎弓屬矢那不中，足令久不得奔，詎教晨光夕昏。\par

\section{相和歌辭 王昭君}
\textbf{崔國輔}\par\vspace{0.5em}
漢使南還盡，胡中妾獨存。\par
紫臺緜望絕，秋草不堪論。\par

\end{document}